\documentclass[12pt,a4paper]{book}
%\documentclass[a4paper,french]{book}
%\usepackage[17pt]{extsizes}


\input{../1ere-preambule}

% Packages nécessaires aux graphiques de l'exercice « Représenter une suite arithmétique »
%\usepackage{tikz}

\newcommand{\va}[1]{\lvert \ #1 \ \rvert}
\newcommand{\vaf}[1]{\left\lvert \ #1\  \right\rvert}

\begin{document}
\graphicspath{{images/}}
%\setcounter{chapter}{1}
\chapter[Suites - Exercices]{Suites arithmétiques \\ Exercices obligatoires}
\thispagestyle{fancy}


\exerciceDS{}

$(u_n)$ est une suite définie par son expression $u_n=5n^2+6$.

\begin{itemMet}
	\item Calculer les termes $u_0$, $u_1$, $u_{10}$ et $u_{50}$.
\end{itemMet}
%$u_0=6$ ; $u_1=11$ ; $u_{10}=506$ ; $u_{50}=12\,506$

\exerciceDS{}

$(v_n)$ est la suite définie par son premier terme $v_0=1$ et la relation de récurrence $v_{n+1}=2v_n+3$.

Pour chaque affirmation, dire si elle est vraie ou fausse en justifiant.
\begin{multicols}{4}
\begin{enumMet}
	\item $v_1=5$.
	\item $v_{10}=2\times v_9+3$.
	\item $v_1=13$.
	\item $v_{15}=2\times 14+3$.
\end{enumMet}
\end{multicols}
%a. Vrai : $v_1=2\times 1+3=5$.
%b. Vrai : c'est la relation de récurrence appliquée à $n=9$.
%c. Faux : $v_1=5$.
%d. Faux : la relation de récurrence fait intervenir $v_{14}$ et non 14.

\exerciceDS{}

$(u_n)$ est la suite définie par son expression $u_n=12+3n$.
\begin{enumMet}
	\item Calculer les termes $u_1$, $u_2$ et $u_{15}$.
	\item Existe-t-il une valeur entière de $n$ telle que $u_n=57$ ? et telle que $u_n=58$ ?
	\item Existe-t-il une valeur entière de $n$ telle que $u_n=58$ ?
\end{enumMet}
%1. $u_1=15$ ; $u_2=18$ ; $u_{15}=57$.
%2. $12+3n=57 \Leftrightarrow n=15$ : oui. $12+3n=58 \Leftrightarrow n=\frac{46}{3}$ : non.

\exerciceDS{}

Pour chacune des suites ci-dessous, trouver les trois termes suivants et conjecturer une formule générale.

\begin{enumMet}
	\item $a_0=3$ ; $a_1=6$ ; $a_2=9$ ; $a_3=12$ ; \dots
	\item $b_0=200$ ; $b_1=100$ ; $b_2=50$ ; $b_3=25$ ; \dots
	\item $c_1=1$ ; $c_2=\dfrac{1}{2}$ ; $c_3=\dfrac{1}{3}$ ; \dots
\end{enumMet}

%a. $a_n=3n$ ; b. $b_n=\frac{200}{2^n}$ ; c. $c_n=\frac{1}{n}$ ; 

\exerciceDS{}

\textbf{QCM} Les premiers termes d'une suite $(w_n)$ sont $w_0=60$, $w_1=30$ et $w_2=15$.

Une relation de récurrence de la suite peut être :
\begin{multicols}{2}
\begin{enumMet}
	\item $w_{n+1}=\dfrac{w_n}{2}$.
	\item $w_{n+1}=w_n-\dfrac{30}{n+1}$.
	\item $w_{n+1}=2w_n-90$.
	\item $w_{n+1}=w_n-30$.
\end{enumMet}
\end{multicols}
%Réponse a.

\exerciceDS{}

Pour chaque suite, trouver les trois termes suivants et conjecturer une relation de récurrence.
\begin{enumMet}
	\item $u_0=17$ ; $u_1=171$ ; $u_2=1\,711$ ; \dots
	\item $v_1=563$ ; $v_2=542$ ; $v_3=521$ ; \dots
\end{enumMet}
%a. $e_{n+1}=10 e_n+1$ ; b. $f_{n+1}=f_n-21$


\exerciceDS{}

$(u_n)$ est une suite arithmétique. On sait que $u_0=-0,1$ et $u_1=0,8$.
\begin{itemMet}
	\item Calculer la raison de cette suite, puis $u_2$ et $u_4$.
\end{itemMet}
%$r=0,9$ ; $u_2=1,7$ ; $u_4=3,5$

\exerciceDS{}

$(v_n)$ est une suite arithmétique. On sait que $v_0=0,52$ et $v_3=1,54$.
\begin{itemMet}
	\item Calculer la raison de cette suite, puis $v_2$.
\end{itemMet}
%$v_3=v_0+3r \Leftrightarrow r=\frac{1,54-0,52}{3}=0,34$ ; $v_2=0,52+2\times 0,34=1,2$

\exerciceDS{}

$(w_n)$ est une suite arithmétique de raison $r=0,2$ et on sait que $w_0=12,5$.
\begin{itemMet}
	\item Calculer $w_1$, $w_4$ et $w_{50}$.
\end{itemMet}
%$w_1=12,7$ ; $w_4=13,3$ ; $w_{50}=12,5+50\times 0,2=22,5$

\exerciceDS{}

$(w_n)$ est une suite arithmétique de premier terme $w_2=2,5$ et de raison $r=-0,25$.
\begin{enumMet}
	\item Calculer les termes $w_3$, $w_4$ et $w_5$.
	\item Donner la relation de récurrence donnant $w_{n+1}$ en fonction de $w_n$.
	\item Donner l'expression de $w_n$ en fonction de $n$.
	\item Calculer $w_{30}$.
\end{enumMet}

\begin{battention}
	Attention, ici le premier terme de la suite $(w_n)$ n'est pas $w_0$.
\end{battention}

%1. $w_3=2,25$ ; $w_4=2$ ; $w_5=1,75$
%2. $w_{n+1}=w_n-0,25$
%3. $w_n=w_2+(n-2)\times(-0,25)=-0,25n+3$
%4. $w_{30}=-0,25\times 30+3=-4,5$

\exerciceDS{}

\textbf{Vrai/Faux}
\begin{itemMet}
	\item $(u_n)$ est une suite arithmétique de raison 3 et de premier terme $u_0=1$.
	\item $(v_n)$ est une suite arithmétique de raison $-\dfrac{3}{4}$ et de premier terme $v_0=-\dfrac{1}{4}$.
	\item $(w_n)$ est une suite arithmétique de raison $1,5$ et de premier terme $w_0=0,2$.
\end{itemMet}

Pour chaque affirmation, dire si elle est vraie ou fausse en justifiant.
\begin{multicols}{2}
	\begin{enumMet}
	\item $u_{25}=100$.
	\item $v_{25}$ est le $25^{\text{e}}$ terme de la suite $(v_n)$.
	\item $w_6$ est supérieur à 9.
	\item $v_{12}<v_{13}$.
	\item $v_{25}<w_{25}<u_{25}$.
\end{enumMet}
\end{multicols}
%a. Faux : $u_{25}=1+25\times 3=76$.
%b. Faux : c'est le $26^{\text{e}}$ terme (le premier est $v_0$).
%c. Vrai : $w_6=0,2+6\times 1,5=9,2>9$.
%d. Faux : la raison est négative donc $(v_n)$ est décroissante.
%e. Vrai : $v_{25}=-19$ ; $w_{25}=37,7$ ; $u_{25}=76$.


\exerciceDS{}

Les situations ci-dessous peuvent-elles être modélisées par une suite arithmétique ? Si oui, en déterminer la raison et le premier terme.
\begin{enumMet}
	\item Le prix d'un produit augmente de 2 \% par an. Son prix initial était $2,15$ \euro.
	\item Un magasin de vêtements fait une promotion sur des t-shirts : le premier acheté est au prix de 12 \euro\ et les suivants au prix de 8 \euro\ chacun.
	\item Pendant les premiers jours d'une épidémie, le nombre de patients qui se rendent dans un cabinet médical augmente quotidiennement de 10 patients. Le premier jour, le cabinet reçoit 35 patients.
	\item Le prix d'un produit augmente de 20 centimes par an. Son prix initial était $2,15$ \euro.
\end{enumMet}
%a. Non (suite géométrique de raison 1,02).
%b. Oui : $u_1=12$ et $r=8$.
%c. Oui : $u_1=35$ et $r=10$.
%d. Oui : $u_0=2,15$ et $r=0,2$.

\exerciceDS{}

Un téléphérique progresse à vitesse constante ; chaque seconde, son altitude augmente de $0,85$ m. La gare de départ est à une altitude de 1\,350 m.
\begin{itemMet}
	\item La durée du trajet étant de 15 min, quelle est l'altitude de la gare d'arrivée ?
\end{itemMet}
%$15$ min $=900$ s ; $1350+900\times 0,85=2115$ m

\exerciceDS{}

Voici quatre représentations graphiques de suites numériques.
\begin{enumMet}
	\item Indiquer celles qui correspondent à des suites arithmétiques.
	\item Pour chacune des suites arithmétiques, indiquer le premier terme et préciser si la raison est positive, négative ou nulle.
\end{enumMet}

\begin{center}
\begin{tabular}{cc}
% ---------- suite (u_n) ----------
\begin{tikzpicture}[scale=0.7]
	\draw[help lines,dashed,gray!60] (0,0) grid (6.5,4.5);
	\draw[->] (0,0) -- (7,0) node[below right] {$n$};
	\draw[->] (0,0) -- (0,5) node[above left] {$u_n$};
	\foreach \x in {1,2,3,4,5} \draw (\x,0.1) -- (\x,-0.1) node[below] {\x};
	\foreach \y in {1,2,3,4} \draw (0.1,\y) -- (-0.1,\y) node[left] {\y};
	\node at (-0.25,-0.3) {0};
	\foreach \x/\y in {0/2,1/2,2/2,3/2,4/2,5/2}
		\node[orange] at (\x,\y) {\textbf{+}};
\end{tikzpicture}
&
% ---------- suite (v_n) ----------
\begin{tikzpicture}[scale=0.7]
	\draw[help lines,dashed,gray!60] (0,0) grid (6.5,4.5);
	\draw[->] (0,0) -- (7,0) node[below right] {$n$};
	\draw[->] (0,0) -- (0,5) node[above left] {$v_n$};
	\foreach \x in {1,2,3,4,5} \draw (\x,0.1) -- (\x,-0.1) node[below] {\x};
	\foreach \y in {1,2,3,4} \draw (0.1,\y) -- (-0.1,\y) node[left] {\y};
	\node at (-0.25,-0.3) {0};
	\foreach \x/\y in {0/4,1/3.25,2/2.6,3/2.15,4/1.75,5/1.5}
		\node[blue] at (\x,\y) {\textbf{+}};
\end{tikzpicture}
\\[0.6cm]
% ---------- suite (w_n) ----------
\begin{tikzpicture}[scale=0.7]
	\draw[help lines,dashed,gray!60] (0,0) grid (6.5,4.5);
	\draw[->] (0,0) -- (7,0) node[below right] {$n$};
	\draw[->] (0,0) -- (0,5) node[above left] {$w_n$};
	\foreach \x in {1,2,3,4,5} \draw (\x,0.1) -- (\x,-0.1) node[below] {\x};
	\foreach \y in {1,2,3,4} \draw (0.1,\y) -- (-0.1,\y) node[left] {\y};
	\node at (-0.25,-0.3) {0};
	\foreach \x/\y in {0/0.15,1/0.6,2/1.3,3/2.2,4/3.8}
		\node[violet] at (\x,\y) {\textbf{+}};
\end{tikzpicture}
&
% ---------- suite (t_n) ----------
\begin{tikzpicture}[scale=0.7]
	\draw[help lines,dashed,gray!60] (0,0) grid (6.5,4.5);
	\draw[->] (0,0) -- (7,0) node[below right] {$n$};
	\draw[->] (0,0) -- (0,5) node[above left] {$t_n$};
	\foreach \x in {1,2,3,4,5} \draw (\x,0.1) -- (\x,-0.1) node[below] {\x};
	\foreach \y in {1,2,3,4} \draw (0.1,\y) -- (-0.1,\y) node[left] {\y};
	\node at (-0.25,-0.3) {0};
	\foreach \x/\y in {0/0.5,1/1.2,2/1.95,3/2.7,4/3.4,5/4.1}
		\node[green!60!black] at (\x,\y) {\textbf{+}};
\end{tikzpicture}
\end{tabular}
\end{center}
%a. $(u_n)$ et $(t_n)$ : les points sont alignés.
%b. $(u_n)$ : $u_0=2$, raison nulle. $(t_n)$ : $t_0=0,5$, raison positive.

\exerciceDS{}

Représenter par un nuage de points les quatre premiers termes des suites $(a_n)$ et $(b_n)$, où $n$ est un nombre entier positif.
\begin{enumMet}
	\item $a_0=-4$ et $a_{n+1}=-2,5+a_n$.
	\item $b_n=1,2+1,5\times n$.
\end{enumMet}
%a. $-4$ ; $-6,5$ ; $-9$ ; $-11,5$
%b. $1,2$ ; $2,7$ ; $4,2$ ; $5,7$

\vspace{0.3cm}
\noindent\textbf{Pour les 2 exercices suivants}

\noindent Pour chacune des suites arithmétiques, déterminer son sens de variation en justifiant.

\exerciceDS{}

\begin{enumMet}
	\item $u_0=1$ et, pour tout nombre entier positif $n$, $u_{n+1}=0,85+u_n$.
	\item Pour tout nombre entier positif $n$, $v_n=-\dfrac{1}{4}\times n+5$.
	\item $w_0=0$ et, pour tout nombre entier positif $n$, $w_{n+1}=-2,71+w_n$.
	\item Pour tout nombre entier positif $n$, $t_n=\dfrac{3-2n}{5}$.
\end{enumMet}
%a. $r=0,85>0$ : croissante.
%b. $r=-\frac{1}{4}<0$ : décroissante.
%c. $r=-2,71<0$ : décroissante.
%d. $t_n=-\frac{2}{5}n+\frac{3}{5}$, $r=-\frac{2}{5}<0$ : décroissante.

\exerciceDS{}

\begin{enumMet}
	\item $t_0=-1,2$ et, pour tout nombre entier positif $n$, $t_{n+1}=t_n$.
	\item Pour tout nombre entier positif $n$, $s_n=5,7\times n+4,3$.
	\item $u_1=-3,4$ et, pour tout nombre entier positif $n$, $u_{n+2}=u_{n+1}-0,3$.
\end{enumMet}
%a. $r=0$ : suite constante.
%b. $r=5,7>0$ : croissante.
%c. $r=-0,3<0$ : décroissante.

\exerciceDS{} \textbf{Chauffage géothermique}

Héléna souhaite installer un chauffage géothermique dans sa maison. Une entreprise spécialisée lui propose une étude. Un forage initial de 100 m doit être réalisé et des échangeurs de chaleur doivent être installés dans la maison. Le coût de cette réalisation est de 35\,000 \euro.

Pour améliorer son rendement, l'entreprise suggère de réaliser un forage plus profond. Chaque décamètre supplémentaire est facturé 550 \euro\ (on rappelle que 1 décamètre se note 1 dam et vaut 10 m). L'entreprise remet à Héléna le tableau ci-dessous qui met en relation les forages supplémentaires et le coût de l'installation.

\begin{center}
\begin{tabular}{|c|c|c|}
	\hline
	\textbf{Profondeur totale} & \textbf{Profondeur} & \textbf{Coût total} \\
	\textbf{du forage} (en m) & \textbf{supplémentaire} (en dam) & \textbf{(en \euro)} \\
	\hline
	$100+0$ & 0 & 35\,000 \\
	\hline
	$100+10$ & 1 & 35\,550 \\
	\hline
	$100+20$ & 2 & 36\,100 \\
	\hline
	$100+30$ & 3 & 36\,650 \\
	\hline
	$100+40$ & 4 & 37\,200 \\
	\hline
\end{tabular}
\end{center}

\noindent\textbf{Problématique} : on souhaite estimer la profondeur maximale avec un budget de 40\,000 \euro.

\begin{enumMet}
	\item Le coût de l'installation est représenté par une suite $(u_n)$ où $n$ désigne le nombre de décamètres supplémentaires du forage.
	\begin{enumMet}
		\item Quelle est la nature de la suite $(u_n)$ ? Donner sa raison et son premier terme.
		\item Justifier la phrase : «\,Le coût de l'installation pour un forage de 260 m de profondeur est représenté par $u_{16}$.\,»
		\item Quel est le coût de l'installation pour un forage de 200 m de profondeur ?
	\end{enumMet}
	\item Que renvoie la fonction en Python ci-dessous pour $n=4$, $n=10$, $n=16$ ? Est-ce cohérent avec les résultats précédents ?

\begin{center}
\begin{minipage}{0.55\linewidth}
\begin{verbatim}
def forage(n):
    u = 35000
    for k in range(n):
        u = u + 550
    return u
\end{verbatim}
\end{minipage}
\end{center}

	\item Répondre à la problématique de deux manières différentes.
\end{enumMet}
%1.a. Suite arithmétique de raison 550 et de premier terme $u_0=35\,000$.
%1.b. $260=100+160$ soit 16 dam supplémentaires.
%1.c. $u_{10}=35\,000+10\times 550=40\,500$ \euro.
%2. $37\,200$ ; $40\,500$ ; $43\,800$.
%3. $35\,000+550n\leqslant 40\,000 \Leftrightarrow n\leqslant 9,09$, soit 9 dam : profondeur maximale 190 m.

\exerciceDS{} \textbf{Politique salariale}

Une entreprise propose à ses cadres administratifs deux types de contrat lors de leur embauche.

\begin{itemMet}
	\item \textbf{Contrat A} : 25\,000 \euro\ la première année, puis 750 \euro\ d'augmentation par an.
	\item \textbf{Contrat B} : 22\,000 \euro\ la première année, puis 1\,000 \euro\ d'augmentation chaque année.
\end{itemMet}

\noindent\textbf{Problématique} : l'objectif est de déterminer le contrat le plus avantageux pour une candidate, selon la durée pendant laquelle elle reste en poste.

On considère $u_n$ et $v_n$ les sommes perçues la $n$-ième année respectivement avec les contrats A et B.

\begin{enumMet}
	\item On souhaite modéliser la situation à l'aide de la feuille de calcul ci-dessous.

\begin{center}
\begin{tabular}{|c|c|c|c|}
	\hline
	 & \textbf{A} & \textbf{B} & \textbf{C} \\
	\hline
	\textbf{1} & $n$ & u(n) & v(n) \\
	\hline
	\textbf{2} & 1 & 25 000 & \\
	\hline
	\textbf{3} & 2 & & \\
	\hline
	\textbf{4} & 3 & & \\
	\hline
	\textbf{5} & 4 & & \\
	\hline
	\textbf{6} & 5 & & \\
	\hline
\end{tabular}
\end{center}

	\begin{enumMet}
		\item Quelle valeur doit être saisie dans la cellule C2 ?
		\item Quelles formules, à recopier vers le bas, peut-on saisir dans les cellules B3 et C3 ?
		\item Utiliser la calculatrice ou le tableur pour répondre à la problématique.
	\end{enumMet}
	\item \begin{enumMet}
		\item Expliquer en quoi la fonction en Python ci-dessous permet également de répondre à la problématique.

\begin{center}
\begin{minipage}{0.65\linewidth}
\begin{verbatim}
def comparaison(u, v):
    n = 1
    while u > v:
        u = u + 750
        v = 22000 + 1000 * n
        n = n + 1
    return n
\end{verbatim}
\end{minipage}
\end{center}

		\item Quelle commande doit-on saisir dans la console pour répondre à la problématique ?
		\item Quel résultat obtient-on ? Est-ce cohérent avec les résultats précédents ?
	\end{enumMet}
\end{enumMet}
%1.a. 22 000
%1.b. B3 : =B2+750 ; C3 : =C2+1000
%1.c. $u_n=25\,000+750(n-1)$ et $v_n=22\,000+1\,000(n-1)$ ; $v_n\geqslant u_n$ à partir de $n=13$.

\end{document}
